%%%

%%%   Самарский государственный технический университет

%%%   Кафедра прикладной математики и информатики

%%%

%%%   Преамбула для статьи в журнал "Вестник Самарского государственного

%%%   технического университета. Серия: «Физико-математические науки»"

%%%   Ориентирована под MikTEx 2.4 for Win32

%%%

%%%   Хотя сам журнал собирается под GNU/Linux на дистрибутиве TeXLive



\documentclass[11pt,a4paper,twoside]{article}



\usepackage[cp1251]{inputenc}

\usepackage[T2A]{fontenc}

\usepackage[english,russian]{babel}

\usepackage{samgtubib}

\usepackage[dvips]{graphicx}

\usepackage[multiple]{footmisc}



\usepackage{samgtu}

\renewcommand{\VestnikYear}{2009} % Год издания Вестника

\renewcommand{\VestnikNumber}{2\,(19)} % Номер Вестника

\renewcommand{\VestnikMonth}{Ноябрь} % Месяц выхода Вестника

\renewcommand{\VestnikData}{01.11.09} % Дата подписания Вестника в печать

\renewcommand{\VestnikUPL}{26,64} % Усл. печ. л.

\renewcommand{\VestnikUIL}{25,73} % Уч.-изд. л.

\renewcommand{\VestnikRegNumber}{479/09} % Регистрационный номер

\renewcommand{\VestnikOrder}{994} % Номер заказа









%%% Attention:

%%% Correct command \Im



%\samgtupubl % для генерации pdf, отдаваемого в типографию СамГТУ

\pdfcrop % обрезка pdf для размещения на math-net.ru

\draft % На корректуре

\nofilllastpage % Последняя страница не заполнена не полностью

%\briefcomm % Статья является кратким сообщением



\vsgtu{1333}



\FirstPage{1} \LastPage{12}





\begin{document}



\udk{519.853.4}

\msc{90C26}



\rutitle{Об одном методе поиска глобального экстремума непрерывной

функции на симплексе} {Об одном методе поиска глобального

экстремума \dots}



\ruauthor{М.~Ю.~Лившиц, А.~П.~Сизиков}{М.~Ю.~Л\,и\,в\,ш\,и\,ц,

А.~П.~С\,и\,з\,и\,к\,о\,в}



\ruaffil{\rusamgtu.} \enaffil{\ensamgtu.}



\ruauthordetails{3}{

    \fullauthor{\rupid{?????}{Михаил Юрьевич Лившиц}} \regalia{д.т.н., проф.;

    \mem{mikhaillivshits@gmail.com}} \position{заведующий кафедрой} \dept{каф. управления и системного анализа теплоэнергетических и социотехнических комплексов}\\

    \fullauthor{\rupid{?????}{Александр Павлович Сизиков}} \regalia{к.э.н., доц.;\mem{apsizikov@mail.ru}; автор, ведущий переписку} \position{докторант}

    \dept{каф. управления и системного анализа теплоэнергетических и социотехнических комплексов}

}



\enauthordetails{3}{

    \fullauthor{\enpid{?????}{Mikhail Yu. Livshits}}   \regalia{Dr. Techn. Sci.; \mem{radch@samgtu.ru}} \position{Head of Dept.} \dept{Dept. of

     management and system analysis of thermal power and socio-technical systems}\\

    \fullauthor{\enpid{?????}{Aleksandr P. Sizikov}} \regalia{Ph.D., Assoc.; \mem{apsizikov@mail.ru}; Corresponding Author}

    \position{Doctoral Student} \dept{Dept. of management and system analysis of thermal power and socio-technical systems}

}



\ruabstract{Рассматривается невыпуклая задача математического

программирования, допустимой областью которой является симплекс.

Для приближенного решения задачи предложен двухэтапный алгоритм.

На первом этапе методом $\Psi$"=преобразования определяется

область глобального оптимума, на втором осуществляется локальная

<<доводка>> решения. Метод $\Psi$"=преобразования модифицирован

с~учетом специфики рассматриваемой задачи. Определение

$\Psi$"=функции осуществляется по результатам статистических

испытаний, реализуемых с помощью генератора случайных чисел,

равномерно распределенных в симплексе. Для уточнения решения

применяется предложенный метод отражения правильных симплексов.

Приведен пример применения разработанного алгоритма для решения

задачи оптимизации компонентного состава углеводородной смеси.}



\rukeywords{оптимизация, невыпуклые задачи, метод $\Psi

$"=преобразования, равномерное распределение в симплексе,

многокомпонентные смеси}



\entitle{Method of Searching for~Global Extremum of a~Continuous Function on a~Simplex

}



\enauthor{M.~Yu.~Livshits, A.~P.~Sizikov}{L\,i\,v\,s\,h\,i\,t\,s~M.~Yu.,

S\,i\,z\,i\,k\,o\,v~A.~P.}





\enabstract{A non-convex problem of mathematical programming is

considered, which permissible region is a simplex. A two-stage

algorithm is proposed for approximate solution of the problem. The

region of global optimum is determined using the $\Psi$-transform

method at the first stage; local <<fine-tuning>> of the solution

is performed at the second stage. The $\Psi$-transform was

modified taking into account the special features of the problem

under consideration. $\Psi$-function is determined according to

the results of statistical tests implemented using the generator

of ran-dom numbers uniformly distributed over the simplex. The

proposed method of reflection of regular simplexes is used for

fine-tuning of the solution. An example of application of the

developed algorithm for solving the problem of optimization of

component composition of the hydrocarbon mixture is presented.}



\enkeywords{optimization, non-convex problems, $\Psi$-transform

method, uniform distribution over the simplex, multi-component

mixtures}



\date{20/IV/2015}

\revisiondate{27/VI/2015} \accepteddate{08/VIII/2015}



\makerutitle



В широкой области технических приложений возникает необходимость решения

нелинейной и, в общем случае, невыпуклой задачи:

\begin{gather}

\label{livshits:eq1}

f(x)\to \max,\\

\label{livshits:eq2}

\begin{cases}

\sum\limits_{j=1}^m x_j =1,\\

x\ge 0,

\end{cases}

\end{gather}

где $f(x)$ "--- некоторая непрерывная функция $x=\big(x_1,x_2 ,\ldots,x_m \big)$.



Для решения невыпуклых оптимизационных задач наиболее часто используют

стохастический подход. Он включает широкий спектр методов "--- от простого

случайного поиска до различных рандомизированных эвристик \cite{livsh:bib1,livsh:bib2}. Все эти

методы так или иначе используют глобальный и локальный поиск в~различных

реализациях и сочетаниях.



В методе имитации отжига переход от глобального к локальному поиску

осуществляется постепенно \cite{livsh:bib3,livsh:bib4}. На~начальном этапе осуществляется случайный

поиск по всей допустимой области. Но~затем, оставаясь случайным, он все

более и более локализуется, становится избирательным. В~методе

кросс"=энтропии [\citen{livsh:bib5}--\citen{livsh:bib7}] идея постепенной локализации поиска реализуется через

механизм изменения распределения поисковой выборки с тем, чтобы

перспективные варианты выпадали с большей вероятностью. Существуют методы,

использующие так называемый <<мультистарт>> локальных поисков

(детерминированных или рандомизированных) сразу из множества точек области

определения \cite{livsh:bib8,livsh:bib9}. Для решения некоторых трудно формализуемых задач

эффективны генетические алгоритмы [\citen{livsh:bib10}--\citen{livsh:bib12}]. В целом же, эти алгоритмы

обладают всеми достоинствами и недостатками рандомизированных эвристик \cite{livsh:bib12}.



Авторы предлагают для решения задачи \eqref{livshits:eq1}, \eqref{livshits:eq2} модификацию известного метода

$\Psi$"=преобразования \cite{livsh:bib14,livsh:bib15}, сочетающего детерминированный и~стохастический подходы. По сравнению с методами \cite{livsh:bib1,livsh:bib12} предложенный метод

менее чувствителен к~размерности задачи. В~его основе лежит разбиение,

применяемое для построения интеграла Лебега. В~рассмотрение вводится

монотонно убывающая функция $\Psi (\zeta )=m(E)$, которая определяется

как мера Лебега некоторой измеримой области $E(\zeta )=\{x\vert x\in

E,\,f(x)\ge \zeta \}$ евклидова пространства. Ноль функции $\Psi (\zeta )$

соответствует значению глобального экстремума целевой функции \eqref{livshits:eq1}.

Определение функции $\Psi (\zeta )$ в~аналитическом виде затруднительно,

поэтому некоторое ее приближение строится по~точкам, которые получают с~помощью статистических испытаний.



В поставленной задаче \eqref{livshits:eq1}, \eqref{livshits:eq2} прежде чем использовать алгоритм $\Psi

$"=преобразования, необходимо при статистических испытаниях решить вопрос

генерирования случайных точек, равномерно распределенных в~области \eqref{livshits:eq2}.



Правильный $(m-1)$"=мерный симплекс \eqref{livshits:eq2}, можно поместить в~соответствующий

$(m-1)$"=мерный единичный гиперкуб $R^{m-1}$, получение равномерного

распределения случайных чисел в котором "--- задача тривиальная. Поскольку

симплекс является подмножеством гиперкуба, точки, попадающие в симплекс,

также равномерно распределены. Однако, для поставленной задачи \eqref{livshits:eq1}, \eqref{livshits:eq2}

этот путь непродуктивен. Как известно, объем правильного $(m-1)$"=мерного

симплекса со стороной единичной длины равен

\[

\frac{1}{(m-1)!}\sqrt {\frac{m}{2^{m-1}}},

\]

и при росте $m$ быстро стремится к нулю. Соответственно, стремится к~нулю и~вероятность того, что случайная точка, сгенерированная в~гиперкубе,

окажется в симплексе. Поэтому для получения случайных допустимых вариантов

используем алгоритм \cite{livsh:bib16,livsh:bib17}, который обеспечивает перенос равномерно

распределенных случайных точек в симплекс \eqref{livshits:eq2} путем аффинных и линейных

преобразований, не влияющих на закон распределения.



Алгоритм 1.



1. Генерируются $(m-1)$ случайных чисел, равномерно распределенных на~единичном отрезке. И~эти числа нумеруются в~порядке возрастания: $\xi _1 \le

\xi _2 \le \ldots \le \xi _{m-1}$.



2. Формируется вектор $\xi'=\left(\xi_1', \xi_2', \ldots, \xi'_{m-1} \right)$,

где $\xi'_1 =\xi_1$, $\xi'_2 =\xi _2 -\xi _1$, {\ldots}, $\xi

'_{m-1} =\xi _{m-1} -\xi _{m-2}$. Этот вектор, как показано в работе \cite{livsh:bib16},

равномерно распределен на множестве

$$

M=\left\{x\in R^{m-1}\vert x_1 \ge 0,x_2 \ge

0,\ldots,x_{m-1} \ge 0,\sum\limits_{i=1}^{m-1} x_i \le 1 \right\}.$$



3. Все угловые точки многогранника $M$, кроме нулевой, лежат на расстоянии

$\sqrt 2$ друг от друга. Многогранник $M$ растягивается до правильного

$\tilde {M}$. Для этого перенесем вершину, лежащую в начале координат так,

чтобы расстояние от нее до всех остальных вершин было также $\sqrt 2$. Это

будет точка с координатами $\left(x_1^0 ,x_2^0 ,\ldots,x_{m-1}^0 \right)$, где $x_i^0

=r=(1-\sqrt m )/(m-1)$, $i=1,2,\ldots,m-1.$ Получился многогранник,

конгруэнтный $(m-1)$"=мерному симплексу в $m$"=мерном пространстве. Обозначим

через $\tilde {\xi}$ точку множества $\tilde {M}$, отображающую точку

$\xi'$ множества $M$. Точку $\tilde {\xi }$ определяем в~результате отражения $\xi'$ относительно гиперплоскости $\left \{x\in

R^{m-1}\vert \sum\limits_{i=1}^{m-1} {x_i =1} \right \}$. Для этого сначала находим

проекцию $\xi'$ на эту гиперплоскость:

$\check{\xi}_i =\xi'_i +\Delta$, $i=1,2,\ldots,m-1$, где $\Delta $ определяем из~условия $\sum\limits_{i=1}^{m-1}

\check{\xi}_i =1$ по формуле $\Delta =\cfrac{1}{m-1}\left( 1-\sum\limits_{i=1}^{m-1}

\xi'_i\right)$. Отражение находим с~учетом того, что

коэффициент растяжения равен $\sqrt m $: $\tilde {\xi

}=\check {\xi} +\sqrt m \left(\xi'-\check{\xi}\right)$. Это аффинное преобразование не оказывает влияние на закон

распределения, поэтому точка $\tilde {\xi }$ имеет равномерное распределение

в $\tilde {M}$.



4. Точка $\tilde {\xi }$ переносится в $(m-1)$"=мерный симплекс $m$"=мерного пространства переменных задачи \eqref{livshits:eq1}, \eqref{livshits:eq2}. Для этого представляем ее

как выпуклую линейную комбинацию угловых точек $\tilde {M}$ $m$"=мерного пространства. Коэффициенты этой комбинации $\alpha _1$, $\alpha_2$, \ldots, $\alpha_m$ есть

не что иное, как новые координаты точки $\alpha$. Они определяются

следующим образом:

\[

\begin{pmatrix}

\alpha _1  \\

\alpha _2  \\

\ldots \\

\alpha _{m-1} \\

\alpha _m

\end{pmatrix}=

\begin{pmatrix}

 1  & 0  & \ldots & 0  & r  \\

 0  & 1  & \ldots & 0  & r  \\

 \ldots & \ldots & \ldots & \ldots & \ldots \\

 0  & 0 & \ldots & 1  & r  \\

 1 & 1  & \ldots & 1  & 1 \\

\end{pmatrix}^{-1}

\begin{pmatrix}

\tilde {\xi }_1 \\

\tilde {\xi }_2  \\

\ldots \\

\tilde {\xi }_{m-1}\\

 1

\end{pmatrix}.

\]

Теперь вернемся к методу $\Psi$"=преобразования.



Алгоритм 2.



1. Выполняется серия случайных проб, число которых $N$ должно быть

достаточным для получения надежных статистических оценок. На каждой пробе

$k=1,2,\ldots,N$ с~помощью алгоритма~1 генерируется допустимое значение

аргумента $x^{(k)}$ и находится $f\big(x^{(k)}\big)$. По результатам всех проб

определяются среднее арифметическое $f_c$ и максимальное $f_m$ наблюдаемых

значений целевой функции:

\[

f_c =\frac{1}{N}\sum\limits_{k=1}^N f\big(x^{(k)}\big),

\quad

f_m =\max \left\{ f\big(x^{(1)}\big),f\big(x^{(2)}\big),\ldots,f\big(x^{(N)}\big) \right\}.

\]

Отрезок $\big[f_c,f_m\big]$ разбивается на $L$ (порядка десяти) уровней

$\zeta =\big(\zeta_1,\zeta_2 ,\ldots,\zeta _L \big)$, где $\zeta _l =f_c +(l-1)\big(f_m

-f_c \big)L^{-1}$.



2. Выполняется вторая серия случайных проб. Для каждого уровня

$l=1,2,\ldots,L$ подсчитывается $n_l$ "--- число проб, в~которых

выполнилось условие $f \big(x^{(k)} \big)\ge \zeta _l$, и средняя координата:

\[

\bar {x}_{il} =\frac{1}{n_l }\sum\limits_{k=1}^N \delta_l x_i^{(k)}

, \quad i=1,2,\ldots, m, \quad \delta _l = \begin{cases}

1, \quad f\big(x^{(k)} \big)\ge \zeta_l, \\

0, \quad f\big(x^{(k)} \big)<\zeta_l.

\end{cases}

\]

По результатам формируются последовательности:

\begin{equation}

\label{livshits:eq3}

\psi =\{\psi _l \}_{l=1}^L ,

\end{equation}

где $\psi _l =n_l N^{-1}$,

\begin{equation}

\label{livshits:eq4}

 \tilde {x}_i =\left \{\bar {x}_{il} \right\}_{l=1}^L,

\quad i=1,2,\ldots,m.

\end{equation}



3. Определяются квадратичные тренды. Пусть

\[

\Lambda =\begin{pmatrix}

 1  & 1  & 1  \\

 1  & 2  & 4  \\

\ldots & \ldots & \ldots \\

 1  & L  & L^2

\end{pmatrix},

\]

тогда параметры $a=\big(a_0,a_1,a_2\big)^T$ тренда $\psi (l)=a_0 +a_1 l+a_2

l^2$ рассчитываются по формуле $a=\big(\Lambda ^T \Lambda \big )^{-1}\Lambda^T \psi

^T$. Аналогично, параметры $b_i =\big( b_{i0}, b_{i1}, b_{i2} \big)^T$ тренда $\tilde

{x}_i (l)=b_{i0} +b_{i1} l+b_{i2} l^2$ "--- по формуле $b_i =\big(\Lambda

^T\Lambda \big)^{-1}\Lambda ^T\tilde {x}_i^T$.



4. Определяется наименьший положительный корень $l^\ast$ уравнения

$\psi (l)=0$. Его существование обусловлено тем, что $\psi (l)$ на

отрезке $[1,L]$ представляет собой тренд последовательности \eqref{livshits:eq3},

монотонно убывающей до значения, близкого к нулю. Подстановкой $l^\ast$ в~тренды $\tilde {x}(l)$ находится приближенное решение задачи: $x_i^\ast

=\tilde {x}_i (l^\ast )$, $i=1,2,\ldots,m$.



Назначение описанного алгоритма 2 "--- войти в область глобального экстремума.

Предположим, что это удалось. Перейдем к локальному поиску самого

экстремума. Методы, предполагающие использование производных первого или

второго порядков (метод сопряженных градиентов, метод Ньютона) здесь не

применимы по причине возможной негладкости целевой функции $f(x)$. В этом

случае необходимо использовать метод, для реализации которого достаточно

значений самой целевой функции $f(x)$. Например, последовательное отражение

симплексов~\cite{livsh:bib18}.



Классический метод последовательного отражения симплексов "--- это метод

поиска экстремума в $m$"=мерном евклидовом пространстве с помощью $m$"=мерных

симплексов~\cite{livsh:bib18}. В задаче \eqref{livshits:eq1}, \eqref{livshits:eq2} речь идет о~поиске экстремума \mbox{в~$(m-1)$"=мерном} симплексе с помощью $(m-1)$"=мерных же симплексов.



Построим в $m$"=мерном евклидовом пространстве правильный $m$"=мерный

симплекс, одна из вершин которого (назовем ее нулевой) лежит в~начале

координат. Пусть длина ребра равна $\Delta x$, тогда координаты вершин

соответствуют столбцам матрицы:

\[

\begin{pmatrix}

x_1^{(0)} & x_1^{(1)} & x_1^{(2)} & \ldots & x_1^{(m)} \\

x_2^{(0)} & x_2^{(1)} & x_2^{(2)} & \ldots & x_2^{(m)}  \\

\ldots  & \ldots & \ldots  & \ldots  & \ldots \\

x_m^{(0)} & x_m^{(1)} & x_m^{(2)}  & \ldots & x_m^{(m)}

\end{pmatrix} =

\begin{pmatrix}

 0  & \hat{q} & q & \ldots & q  \\

 0  & q  & \hat{q}  & \ldots & q \\

 \ldots & \ldots & \ldots & \ldots & \ldots \\

 0  & q  & q  &  \ldots  & \hat{q}

\end{pmatrix},

\]

где $\hat{q} =\cfrac{\Delta x}{m\sqrt 2 }\left( \sqrt {m+1} +m-1 \right)$,

$q=\cfrac{\Delta x}{m\sqrt 2 }\left( \sqrt {m+1} -1 \right)$.



Перенесем этот симплекс из начала координат так, чтобы все его вершины,

кроме $x^{(0)}$ расположились на гиперплоскости $x_1 +x_2 +\ldots+x_m =1$

и чтобы $(m-1)$"=мерный симплекс с вершинами $x^{(1)}$, $x^{(2)}$, \ldots, $x^{(m)}$

оказался в области начала поиска.



При параллельном переносе нулевой вершины в точку $\tilde {x}^{(0)}$

координаты остальных вершин становятся равными $\tilde {x}^{(k)}=\tilde

{x}^{(0)}+x^{(k)}$, $k=1,2,\ldots,m$. Найдем $\tilde {x}^{(0)}$ из условия

$\tilde {x}_1^{(k)} +\tilde {x}_2^{(k)} +\ldots+ \tilde {x}_m^{(k)} =1$.

Достаточно взять какую"=нибудь одну вершину, например, первую

\begin{multline*}

\sum\limits_{i=1}^m \tilde {x}_i^{(1)} =

\sum\limits_{i=1}^m \tilde {x}_i^{(0)}

+\frac{\Delta x}{m\sqrt 2 }\left( \sqrt {m+1} +m-1 \right)

+\sum\limits_{i=2}^m \frac{\Delta x}{m\sqrt 2 }\left( \sqrt {m+1} -1

\right) =\\

= \sum\limits_{i=1}^m \tilde {x}_i^{(0)}  +\frac{\Delta x}{m\sqrt 2

}\left( \sqrt {m+1} +m-1+\sum\limits_{i=1}^{m-1} \left( \sqrt {m+1} -1

\right) \right)=\\

= \sum\limits_{i=1}^m \tilde {x}_i^{(0)} +\frac{\Delta x}{m\sqrt 2

}m\sqrt {m+1}= 1.

\end{multline*}



Отсюда получаем:

\[

\sum\limits_{i=1}^m \tilde {x}_i^{(0)} =

1-\frac{\Delta x}{\sqrt 2 }\sqrt {m+1}.

\]



Определим начальное расположение $m$"=мерного симплекса, привязав его к~точке

$x^\ast$, полученной методом $\Psi$"=преобразования. В~этом случае

координаты нулевой вершины симплекса определятся следующим образом:

\[

\tilde {x}_i^{(0)} =

x_i^\ast \left( 1-\frac{\Delta x}{\sqrt 2 }\sqrt {m+1} \right),

\quad

i=1,2,\ldots,m.

\]



Координаты остальных вершин, которые и образуют начальный <<рабочий>>

симплекс, представлены столбцами матрицы:

\[

\begin{pmatrix}

 \tilde {x}_1^{(1)} & \tilde {x}_1^{(2)} & \ldots & \tilde {x}_1^{(m)} \\

 \tilde {x}_2^{(1)} & \tilde {x}_2^{(2)} & \ldots & \tilde {x}_2^{(m)} \\

 \ldots & \ldots & \ldots & \ldots \\

 \tilde {x}_m^{(1)} & \tilde {x}_m^{(2)} & \ldots & \tilde {x}_m^{(m)} \\

\end{pmatrix},

\]

где

\[

\tilde {x}_i^{(k)} =

x_i^\ast \left( 1-\frac{\Delta x}{\sqrt 2 }\sqrt {m+1} \right)+

\begin{cases}

\cfrac{\Delta x}{m\sqrt 2 }\left( \sqrt {m+1} +m-1 \right), \quad i=k, \\

\cfrac{\Delta x}{m\sqrt 2 }\left( \sqrt {m+1} -1 \right), \quad i\ne k.

\end{cases}

\]



Процесс поиска состоит в~реализации последовательности шагов, на каждом из

которых происходит построение нового рабочего симплекса путем отражения

вершины с минимальным значением целевой функции $f(x)$ от

противоположной грани исходного симплекса. Пусть $v\in [1,m]$ "--- номер

такой вершины, тогда отраженная вершина определяется по формуле:

\[

\check{x}^{(v)}=\frac{2}{m-1}\left( \sum\limits_{k=1}^m x^{(k)} -x^{(v)} \right)-x^{(v)}.

\]



Покажем, что

$\check{x}^{(v)}$ лежит на гиперплоскости $x_1 +x_2 +\ldots+x_m =1$. Действительно,

\begin{multline*}

\sum\limits_{i=1}^m

\check {x}_i^{(v)} = \frac{2}{m-1}\sum\limits_{i=1}^m \left( \sum\limits_{k=1}^m

x_i^{(k)} -x_i^{(v)}  \right) - \sum\limits_{i=1}^m x_{i=1}^{(v)} = \\

=\frac{2}{m-1}\left( \sum\limits_{k=1}^m \sum\limits_{i=1}^m x_i^{(k)}

-\sum\limits_{i=1}^m x_i^{(v)} \right)-\sum\limits_{i=1}^m

x_i^{(v)} =

\frac{2}{m-1}( m-1)-1=1.

\end{multline*}



Поскольку поиск проводится в ограниченном пространстве, может оказаться, что

$\check {x}^{(v)}$ выходит за допустимую область. Предположим для простоты, что на

текущем шаге алгоритма "--- $f\left( x^{(1)}\right)<f\left( x^{(2)}\right)<\ldots<f\left( x^{(m)} \right)$. Тогда нужно последовательно

перебирать вершины до тех пор, пока не получим допустимого отражения.

Отражение лучшей точки будет признаком окончания текущей серии. Следующая

серия начинается с уменьшения длины ребра $\Delta x$ и переноса нового

<<рабочего>> симплекса в область лучшей точки предыдущей серии.



Описанная методика может быть использована для поиска оптимального

компонентного состава различных смесей, в~частности, для расчета рецептур

смешения товарных нефтепродуктов. Некоторые показатели качества смесей

рассчитываются по специальным, иногда довольно сложным эмпирическим формулам

\cite{livsh:bib19,livsh:bib20}.



Пусть $Q$ "--- множество индексов контролируемых показателей некоторой смеси, и~пусть $\varphi _q (x)$ "--- зависимость (линейная или нелинейная) $q$-го

показателя смеси от компонентов смешения, взятых в пропорциях $x=(x_1 ,x_2

,\ldots,x_m )$. Тогда задачу смешения можно сформулировать так:

\begin{gather*}

C=\sum\limits_{j=1}^m c_j x_j  \to \min ,\\

\begin{cases}

 \phi_q (x)\in \left[p_q^-,p_q^+ \right],q\in Q, \\

 \sum\limits_{j=1}^m x_j =1,\\

 x\ge 0,

\end{cases}

\end{gather*}

где $c_j$ "--- цена $j$-го компонента, $p_q^-$, $p_q^+ $ "--- граничные

значения для $q$-го показателя.



Требования по качеству смеси удобно учесть в виде штрафной составляющей

целевой функции:

\begin{gather}

\label{livshits:eq5}

 f(x)=\sum\limits_{j=1}^m c_j x_j +\mu

\sum\limits_{q\in Q} \varphi _q^2 (x) \to \min,\\

\label{livshits:eq6}

\begin{cases}

\sum\limits_{j=1}^m x_j =1, \\

x\ge 0,

\end{cases}

\end{gather}

где $\mu $ "--- весовой коэффициент штрафной составляющей; $\varphi _q(x)$ "--- процентное отклонение расчетного значения $q$-го показателя от границы

заданного интервала, определяемое по формуле:

\[

\varphi _q (x)=\begin{cases}

10^2\left( p_q^- -\phi _q (x) \right)/p_q^-, \quad p_q^- -\phi _q (x)>0, \\

10^2\left( \phi _q (x)-p_q^+ \right)/p_q^+ , \quad p_q^+ -\phi _q (x)<0, \\

0, \quad \phi _q (x)\in \left [p_q^-,p_q^+  \right].

\end{cases}

\]



Для примера решена задача составления трехкомпонентной смеси со~следующими

условиями. Стоимость $C=1,5x_1 +1,3x_2 +1,0x_3$. Требования к~показателям

(октановому числу, плотности, содержанию серы):

\begin{equation}

\label{livshits:eq7}

\begin{cases}

\phi_1 (x)=90,2x_1 +88,5x_2 +73,6x_3 -14,3x_1 x_2 -9,8x_1 x_3 +28,4x_1

x_2 x_3 \ge 85, \\

\phi _2 (x)=0,82x_1 +0,59x_2 +0,67x_3 \le 0,68, \\

\phi _3 (x)=0,001x_1 +0,003x_2 +0,002x_3 \le 0,002.

\end{cases}

\end{equation}



При составлении функции $f(x)$ в \eqref{livshits:eq5} используется весовой коэффициент $\mu

=10$. По результатам предварительного статистического исследования области

изменения $f(x)$ сформировано 10 уровней с шагом $\Delta \zeta =-70$,

начиная с $\zeta _1 =700$.



Для каждого из них выполнено 1000 статистических испытаний. Получены

последовательности \eqref{livshits:eq3}, \eqref{livshits:eq4} и построены тренды:



a) $\psi (l)=0,6916-0,02906l-0,00374l^2$,



b) $\tilde {x}_1 (l)=0,33041-0,00145l+0,000098l^2$,



c)$\tilde {x}_2 (l)=0,327667+0,013837l-0,00072l^2$,



d) $\tilde {x}_3 (l)=0,341923-0,012392l+0,000626l^2$.





\begin{figure}[h!]\centering \small

\includegraphics[scale=0.9]{livshits1.eps} \par {\sl а \hspace{70mm} б}\\

\vspace{5mm}

\includegraphics[scale=0.9]{livshits2.eps} \par {\sl в \hspace{70mm} г}

\caption*{Результаты статистических испытаний: 1 "---

последовательности (\ref{livshits:eq3}), (\ref{livshits:eq4}); 2

"--- тренды последовательностей (Results of statistical tests: 1

"--- sequences (\ref{livshits:eq3}), (\ref{livshits:eq4});  2 "---

trends of sequences)}

\end{figure}



Из уравнения $\psi (l)=a_0 +a_1 l+a_2 l^2=0$ получено $l^\ast =10,2554$.

Подстановкой $l^\ast $ в тренды $\tilde {x}_1 (l)$, $\tilde {x}_2 (l)$,

$\tilde {x}_3 (l)$. найдено приближенное решение: $x_1^\ast =0,3259$,

$x_2^\ast =0,3934$, $x_3^\ast =0,2807$. Значение целевой функции

\eqref{livshits:eq5}: $f(x^\ast )=115,3$.



Уточнение решения осуществляется по описанному выше модифицированному методу

отражения правильных симплексов. Проведено 20 серий с уменьшением длины

ребра симплекса по формуле $\Delta x_n =0,1n^{-1}$, где $n$ "--- номер

серии. Результаты первого приближения $x^\ast$ и уточненного

$\check{x}$ решения представлены в таблице.



\begin{table}[h!]\centering \small

\begin{tabular}{|c|c|c|c|c|c|c|c|} \hline

 &

$x_1 $&

$x_2 $&

$x_3 $&

$\phi _1 $&

$\phi _2 $&

$\phi _3 $&

$f(x)$ \\

\hline

$x^\ast $&

0,3259&

0,3934&

0,2807&

83,16&

0,687&

0,00209&

115,3 \\

\hline

$\check{x}$&

0,3702&

0,3990&

0,2308&

83,71&

0,693&

0,00206&

87,6 \\

\hline

\end{tabular}

\end{table}



В результате реализации алгоритма выявлена несовместность системы

ограничений \eqref{livshits:eq7} и получено решение по критерию \eqref{livshits:eq5}, основная составляющая

которого "--- взвешенная сумма квадратов процентных отклонений показателей от

требуемых значений.















\thanks{{\bf ORCIDs}



{\bf Mikhail Yu. Livshits:}

\url{http://orcid.org/0000-0001-8027-2273}



{\bf Aleksandr P. Sizikov:}

\url{http://orcid.org/0000-0001-5067-0676}





}











%\EngRef











\begin{thebibliography}{99}



\Bibitem{livsh:bib1}

\by Zhigljavsky~A., Zilinskas~A.

\book Stochastic Global Optimization

\publ Springer

Science-Business Media

\yr 2008

\totalpages 262



\RBibitem{livsh:bib2}

\by Пантелеев~А.~В.

\book Метаэвристические алгоритмы поиска глобального экстремума

\publaddr  М.

\publ МАИ-ПРИНТ

\yr 2009

\totalpages 159

\misctransl

\lang In Russian

\by Panteleev~A.~V.

\book Metaheuristic algorithms for global extremum searching

\publaddr M.

\publ MAI-PRINT publishing house

\yr 2009

\totalpages 159





\Bibitem{livsh:bib3}

\by Jones M.~T.

\book AI Application Programming

\publaddr M.

\publ DMK Press

\yr 2004

\totalpages 312



\RBibitem{livsh:bib4}

\by Савин~А.~Н., Тимофеева~Н.~Е.

\paper Применение алгоритма оптимизации методом

имитации отжига на системах параллельных и распределённых вычислений

\jour Известия Саратовского университета. Новая серия. Серия Математика. Механика.

Информатика

%Саратов: Изд-во Сарат. ун-та,

\yr 2012

\vol 12

\issue 1

\pages 110--116

\misctransl

\lang In Russian

\by Savin~A.~N., Timofeeva~N.~E.

\paper Application of optimization algorithm using the method of simulated annealing for systems of parallel

and distributed computing

\jour Transactions of Saratov University. The new series. Series Mathematics. Mechanics. Informatics

\yr 2012

\vol 12

\issue 1

\pages 110--116



\Bibitem{livsh:bib5}

\by  Botev~Z.~I.,  Kroese~D.~P.

\paper Global likelihood optimization via the

crossentropy method with an application to mixture models

\inbook In Proceedings of

the 36th Winter Simulation Conference

\pages 529--535

\publaddr Washington, D.C.

\yr 2004



\Bibitem{livsh:bib6}

\by Ernst~D., Glavic~M., Stan~G.-B.,  Mannor~S.,  Wehenkel~L., Gif sur

Yvette Supelec

\paper  The cross-entropy method for power system combinatorial

optimization problems

\inbook In Power Tech, 2007 IEEE Lausanne

\pages 1290--1295

\publaddr Lausanne

\yr 2007



\Bibitem{livsh:bib7}

\by Evans~G.~E., Keith~J. M., Kroese~D.~P.

\paper Parallel cross-entropy

optimization

\inbook In Proceedings of the 2007 Winter Simulation Conference

\pages 2196--2202

\publaddr Washington, D.C.

\yr 2007



\RBibitem{livsh:bib8}

\by  Жиглявский~А.~А., Жилинская~А.~Г.

\book Методы поиска глобального экстремума

\publaddr М.

\publ Наука

\yr 1991

\totalpages  248

\misctransl

\lang In Russian

\by Zhigliavsky~A.~A., Zhilinskaya~A.~G.

\book Methods of search for global extremum

\publaddr M.

\publ Nauka

\yr 1991

\totalpages 248



\RBibitem{livsh:bib9}

\by Феоктистов~А.~Г., Горский~С.~А.

\paper  Реализация метода мультистарта в пакете

Градиент

\jour Вестник НГУ Серия: Информационные технологии

\yr 2007

\vol 05

\issue 2

\pages 78--82

\misctransl

\lang In Russian

\by Feoktistov~A.~G., Gorsky~S.~A.

\paper Implementation of multistart method in Gradient package

\jour Vestnik NGU, Series: Information technologies

\yr 2007

\vol 5

\issue 2

\pages 78--82





\Bibitem{livsh:bib10}

\by Cohoon~J.~P., Karro~J., Lienig~J.

\paper Evolutionary Algorithms for the

Physical Design of VLSI Circuits

\inbook  Advances in Evolutionary Computing: Theory

and Applications

\ed  Ghosh~A., Tsutsui~S.

\publ Springer Verlag

\publ London

\yr 2003

\pages 683--712



\RBibitem{livsh:bib11}

\by Гладков~Л.~А., Курейчик~В.~В., Курейчик~В.~М.

\book Генетические алгоритмы

\ed В.~М.~Курейчика.  2-е изд., испр. и доп

\publaddr М.

\publ  Физматлит

\yr 2006

\totalpages 320

\misctransl

\lang In Russian

\by Gladkov~L.~A., Kureychik~V.~V., Kureychik~V.~M.

\book Genetic algorithms

\ed V.~M.~Kureychik. - 2nd edition, revised and enlarged

\publaddr M.

\publ Fizmatlit

\yr 2006

\totalpages 320



\RBibitem{livsh:bib12}

\by Гладков~Л.~А., Гладкова~Н.~В.

\paper Особенности использования нечетких

генетических алгоритмов для решения задач оптимизации и управления

\jour Известия ЮФУ. Технические науки

\yr  2009

\issue 4(93)

\pages 130--136

\misctransl

\lang In Russian

\by Gladkov~L.~A., Gladkova~N.~V.

\paper Special features of applying fuzzy genetic algorithms for solving optimization and control problems

\jour Transactions of YuFU. Engineering sciences

\yr 2009

\issue 4(93)

\pages 130--136





\Bibitem{livsh:bib13}

\by Steven~S. Skiena

\book The Algorithm Design Manual. Second Edition

\publaddr Springer

\yr 2008

\totalpages 730



\RBibitem{livsh:bib14}

\by  Чичинадзе~В.~К.

\book Решение невыпуклых нелинейных задач оптимизации. Метод

$\Psi$"=преобразования

\publaddr М.

\publ Наука

\yr 1983

\totalpages 256

\misctransl

\lang In Russian

\by Chichinadze~V.~K.

\book Solution of nonconvex nonlinear optimization problems. $\Psi$-transform method

\publaddr M.

\publ Nauka. Main editorial board of physical and mathematical literature

\yr 1983

\totalpages 256





\RBibitem{livsh:bib15}

\by Ахмадиев~Ф.~Г., Гильфанов~Р.~М.

\paper Математическое моделирование и оптимизация

<<состав --- свойство>> многокомпонентных смесей

\jour Математическое моделирование,

численные методы и комплексы программ (в строительстве). Известия КГАСУ

\yr 2012

\issue 2(20)

\misctransl

\lang In Russian

\by Akhmadiev~F.~G., Gil'fanov~R.~M.

\paper Mathematical modeling and "composition-property" optimization of multi-component mixtures

\jour Mathematical modeling, numerical methods and software packages (in construction). Transactions of KGASU

\yr 2012

\issue 2(20)



\RBibitem{livsh:bib16}

\by Новоселов~А.~А.

\paper Равномерное распределение на стандартном симплексе.

[Электронный ресурс]

\inbook Электронная библиотека Попечительского совета

механико"=математического факультета Московского государственного

университета [Офиц. сайт]. URL: http://lib.mexmat.ru/books/31728

\misctransl

\lang In Russian

\by Novoselov~A.~A.

\paper Uniform distribution on the standard simplex.  [Electronic resource]

\inbook Electronic library of the Board of trustees of mechanics and mathematics

faculty of Moscow state University [official website]. URL: http://lib.mexmat.ru/books/31728





\RBibitem{livsh:bib17}

\by   Русин~Ю.~В.

\book Алгоритмы статистического моделирования вероятностных

распределений: текст лекций

\publaddr Ярославль

\publ ЯрГУ

\yr 2006

\misctransl

\lang In Russian

\by Rusin~Yu.~V.

\book Algorithms of statistical modeling of probability distributions: text of lectures

\publ YarGU

\publaddr Yaroslavl

\yr 2006





\RBibitem{livsh:bib18}

\by Зедгинидзе~И.~Г.

\book Планирование эксперимента для исследования

многокомпонентных систем

\publaddr М.

\publ Наука

\yr 1976

\totalpages 390

\misctransl

\lang In Russian

\by Zendginidze~I.~G.

\book Experiment planning for investigation of multicomponent systems

\publaddr M.

\publ Nauka

\yr 1976

\totalpages 390



\RBibitem{livsh:bib19}

\by Чинакал~В.~О.

\paper Оптимизация рецептуры светлых нефтепродуктов

\inbook Оптимизация, исследование операций, бионика

\publaddr М.

\publ Наука

\yr 1973

\pages 198--205

\misctransl

\lang In Russian

\by Chinakal~V.~O.

\paper Composition optimization of light oil products

\inbook Optimization, operations research, bionics

\publaddr M.

\publ Nauka

\yr 1973

\pages 198--205



\RBibitem{livsh:bib20}

\by Никитин~В.~А., Мусаев~А.~А.

\paper  Оптимизация компаундирования углеводородных

смесей

\jour Труды СПИИРАН

\yr 2007

\issue 4

\pages 327--336

\misctransl

\lang In Russian

\by Nikitin~V.~A., Musaev~A.~A.

\paper Optimization of hydrocarbon mixtures compounding

\jour Proceedings of SPIIRAN

\yr 2007

\issue 4

\pages 327--336





\end{thebibliography}



%\EngDate



\makeentitle



%\newpage

%\Russian



%

\end{document}

